\chapter*{Etterord: til dei som skal rive}
\addcontentsline{toc}{chapter}{Etterord: til dei som skal rive}
\markboth{Etterord}{}

Eg har kalla denne boka ei avrekning, og avrekningar endar med eit oppgjer. Men oppgjeret eg ber om er ikkje med personane i faget. Det er med strukturen dei arbeider i. Lat meg seie det reint til kvar av dei som har lese så langt.

\emph{Til læraren.} Du veit alt det eg har skrive; du har visst det lenge, i den uroa du kjenner kvar gong du feller ein dom du ikkje heilt kan grunngje, kvar gong ein student spør «kvifor er denne betre?» og du høyrer deg sjølv svare med ord som ikkje held. Den uroa er ikkje di personlege utilstrekkelegheit. Ho er faget sitt manglande fundament, kjend på kroppen. Du kan halde fram med å undertrykkje henne, eller du kan la henne bli kravet ho eigentleg er: kravet om ein grunn å stå på. Krev han. Du er den einaste som kan, frå innsida.

\emph{Til studenten.} Pensumet er under bygging, og det du krev i studiet formar kva faget vert. Ikkje godta dommar du ikkje får grunngjevne. Krev å få vite kva aksar eit arbeid vart vurdert på, kva som vart vege mot kva. Krev tid i substratet og konsultasjon med nabofaga. Når nokon fortel deg at design er ein eigen kunnskapsform som ikkje treng grunngje seg, høyr kva som faktisk vert sagt: at du skal betale for ei utdanning i ein kompetanse ingen vil definere for deg. Du har rett til å krevje grunnen. Faget skuldar deg han.

\emph{Til akademiet.} Det som manglar er ikkje lenger teori; teorien finst, falsifiserbar, lagd fram. Det som manglar er institusjonell vilje til å la han endre noko. Det er ein hard ting eg ber om, fordi å ta fundamentet inn over seg ville tvinge faget til å slutte å gjere mykje av det faget gjer i dag, og heile sjølvbedraget i kapittel \textsc{vii} er bygd for å sleppe nett det. Men valet er ikkje om faget skal endrast. Materialet og maskinene avgjer det. Valet er om endringa skjer innanfrå, med ein grunn på plass, eller utanfrå, med faget redusert til teikna det sat att med.

Eg skreiv ikkje denne boka for å såre eit fag eg er glad i. Eg skreiv henne fordi eg er glad i det, og fordi det å late som om holet ikkje finst er den einaste sikre måten å la faget bli avleggar på. Ein disiplin som tør sjå sin eigen mangel i auga, kan byggje noko av han. Ein som gøymer han bak prisar og prosessar og prat om vrange problem, kan berre vente på at verda går frå han.

So riv. Ikkje av forakt — av omhug for det som kan stå att når sløret er borte. Under sløret er det ingenting. Men ved sida av holet ligg ein grunn, ferdig lagd, og ventar på eit fag som tør gå frå det eine til det andre. Det finst inga mildare sanning å gje deg. Forma lyg ikkje, og det gjer ikkje fråvêret av henne heller.
